\documentclass[a4paper]{article}

% Basic packages
% Español (México) con XeLaTeX: fontspec para fuentes Unicode, polyglossia
% para la división silábica y los nombres del documento en la variante mexicana.
\usepackage{fontspec}
\usepackage{polyglossia}
\setdefaultlanguage[variant=mexican]{spanish}
% Los nombres chinos van como «pinyin (original)»: xeCJK da a los caracteres
% Han su propia fuente; sin ella XeLaTeX los omite con un aviso.
% FandolSong no cubre algunos caracteres de nombres (U+73FA); WenQuanYi Zen Hei sí.
\usepackage[AutoFallBack=true]{xeCJK}
\setCJKmainfont{FandolSong-Regular.otf}
\setCJKfallbackfamilyfont{\CJKrmdefault}{WenQuanYi Zen Hei}
% polyglossia no traduce los nombres de listings.
\AtBeginDocument{\providecommand{\lstlistingname}{}\renewcommand{\lstlistingname}{Listado}%
  \providecommand{\lstlistlistingname}{}\renewcommand{\lstlistlistingname}{Índice de listados}}
\usepackage{amsmath, amssymb}
\usepackage{graphicx}
\usepackage[margin=2.5cm]{geometry}
\usepackage[most]{tcolorbox}
\usepackage{etoolbox}
\usepackage{listings}
\usepackage{hyperref}
\usepackage{booktabs}
\usepackage{subcaption}
\usepackage{float}

% Highlight boxes
\newtcolorbox{knowledgebox}[1]{
    enhanced, colback=blue!5!white, colframe=blue!75!black, colbacktitle=blue!75!black,
    coltitle=white, fonttitle=\bfseries, title={#1}, attach boxed title to top left={yshift=-2mm, xshift=2mm},
    boxrule=1pt, sharp corners
}

\newtcolorbox{importantbox}[1]{
    enhanced, colback=yellow!10!white, colframe=yellow!80!black, colbacktitle=yellow!80!black,
    coltitle=black, fonttitle=\bfseries, title={#1}, sharp corners
}

\newtcolorbox{warningbox}[1]{
    enhanced, colback=red!5!white, colframe=red!75!black, colbacktitle=red!75!black,
    coltitle=white, fonttitle=\bfseries, title={#1}, sharp corners
}

% Listings
\lstset{
    language=Python,
    basicstyle=\ttfamily\small,
    keywordstyle=\color{blue},
    stringstyle=\color{red!60!black},
    commentstyle=\color{green!60!black},
    breaklines=true,
    frame=single,
    numbers=left,
    numberstyle=\tiny\color{gray},
    captionpos=b,
    extendedchars=true
}

% Front-page metadata
\newcommand{\notetitle}{Superpowers Skills: la biblia del desarrollo de software en la era del Vibe Coding}
\newcommand{\noteauthors}{Elaboradas a partir de materiales públicos del curso}
\newcommand{\notedate}{2025}
\newcommand{\videochannel}{Wudaokou Nash (五道口纳什)}
\newcommand{\videopublishdate}{2025}
\newcommand{\videoduration}{aproximadamente 40 minutos}
\newcommand{\videourl}{}
\newcommand{\videocoverpath}{cover.jpg}
\newcommand{\repourl}{https://github.com/hqhq1025/ai-course-notes}


% Footer with repo info
\usepackage{fancyhdr}
\pagestyle{fancy}
\fancyhf{}
\fancyfoot[C]{\thepage}
\fancyfoot[R]{\footnotesize\color{gray}\href{\repourl}{AI Course Notes}}
\renewcommand{\headrulewidth}{0pt}
\renewcommand{\footrulewidth}{0pt}

\begin{document}

\begin{titlepage}
\centering
{\Large Notas del curso\par}
\vspace{1.2cm}
{\huge\bfseries \notetitle\par}
\vspace{0.8cm}
{\large \noteauthors\par}
\vspace{0.3cm}
{\large \notedate\par}
\vspace{1.2cm}

\ifdefempty{\videocoverpath}{
{\small Al generar las notas, indique la ruta local de la portada del video.\par}
}{
\includegraphics[width=0.82\textwidth,height=0.45\textheight,keepaspectratio]{\videocoverpath}\par
}

\vfill
\begin{tcolorbox}[width=0.9\textwidth, colback=black!2!white, colframe=black!60, sharp corners]
\textbf{Autor o canal del video}: \videochannel\par
\textbf{Fecha de publicación}: \videopublishdate\par
\textbf{Duración del video}: \videoduration\par
\ifdefempty{\videourl}{}{
\textbf{Enlace del video}: \href{\videourl}{\nolinkurl{\videourl}}\par
}
\par
\textbf{Página del proyecto}: \href{\repourl}{\nolinkurl{\repourl}}\par
\end{tcolorbox}
\end{titlepage}

\tableofcontents
\newpage

%% --- Inicio del contenido --- %%

\section{Introducción: la crisis de calidad en la era del Vibe Coding}

En la actualidad, con la rápida popularización de la programación asistida por IA (Vibe Coding), los desarrolladores enfrentan un problema central: \textbf{cómo garantizar la exactitud y la mantenibilidad del código al usar herramientas de programación en lenguaje natural puro como Codex y Claude Code}. Esta entrega gira en torno al proyecto de código abierto \textbf{Superpowers} (más de 70,000 estrellas en GitHub), un conjunto de Skills de alta calidad, para responder sistemáticamente a esta pregunta.

\begin{knowledgebox}{El espectro de herramientas del Vibe Coding}
El docente divide las herramientas de programación con IA actuales en tres niveles:
\begin{enumerate}
    \item \textbf{VS Code (sin complementos de IA)}: un editor de código / IDE puro.
    \item \textbf{Cursor}: asistencia de IA ligera con un modelo Tab (entrenado con RL on-policy a gran escala), que admite desarrollo conversacional e interfaz gráfica de diff.
    \item \textbf{Codex / Claude Code}: CLI de programación en lenguaje natural puro, sin IDE, donde solo se necesita describir el requerimiento y verificar la salida.
\end{enumerate}
\end{knowledgebox}

En los métodos tradicionales, la calidad del código se garantiza mediante procesos como \textbf{TDD} (desarrollo guiado por pruebas), \textbf{Code Review} y \textbf{Pre-commit Hook} (por ejemplo, el análisis estático con Ruff). En la era del Vibe Coding, Superpowers «codifica de manera rígida» estas prácticas de ingeniería de software de primer nivel como el sistema operativo normativo del Agente de IA Coding: mediante instrucciones de Skill, y no scripts de Hook, obliga al Agente a superar las tendencias comunes de los modelos grandes a la alucinación, la pereza y la complacencia.

\begin{importantbox}{Idea central: «lento en realidad es rápido»}
Sin ningún Skill ni ninguna restricción, el modelo puede producir código muy rápido, pero el código es inexacto y no es mantenible. Con las restricciones de Superpowers, aunque cada paso es más cuidadoso, la calidad del código está garantizada y la eficiencia de la iteración a largo plazo es mayor: esto es «lento en realidad es rápido».
\end{importantbox}


\section{Mecanismo de carga de los Skills}

\subsection{Arquitectura del System Prompt en Codex}

Antes de que el usuario emita su primera solicitud, Codex ya ha insertado cuatro mensajes de sistema:

\begin{enumerate}
    \item \textbf{Prompt a nivel System}: restricciones a nivel de sistema.
    \item \textbf{Mensaje Developer}: reglas de permisos y sandbox.
    \item \textbf{Mensaje User (inyectado)}: contiene \texttt{agents.md} y la descripción de todos los Skills disponibles; esta es la \textbf{primera exposición} de los Skills.
    \item \textbf{Mensaje Developer}: define el modo de colaboración (Plan mode / Agent mode).
\end{enumerate}

\subsection{Proceso de emparejamiento y carga de un Skill}

\begin{enumerate}
    \item El usuario emite una solicitud real.
    \item El Coding Agent \textbf{empareja el Skill más adecuado} a partir del query.
    \item El Agent invoca la herramienta \texttt{execute\_commands} y ejecuta el comando \texttt{cat} para cargar el contenido completo del \texttt{skill.md} correspondiente.
    \item El contenido del Skill entra al contexto, y todas las interacciones posteriores se ejecutan a partir de la descripción de ese Skill.
    \item Al terminar un turn, el usuario puede seguir iterando sobre sus necesidades.
\end{enumerate}

\begin{knowledgebox}{Los dos orígenes de agents.md}
\begin{itemize}
    \item \textbf{A nivel de sistema}: \texttt{\~{}/.codex/agents.md} (configuración global).
    \item \textbf{A nivel de proyecto}: el \texttt{agents.md} en la raíz del repositorio de código (instrucciones propias del proyecto).
\end{itemize}
Ambos se inyectan en el tercer mensaje User.
\end{knowledgebox}

\subsection{Instalación y actualización}

\begin{itemize}
    \item Claude Code instala Superpowers mediante el mecanismo de plugin.
    \item Codex deja que el Agent se instale a sí mismo a partir de lenguaje natural.
    \item Forma de actualizar: entrar al directorio de Superpowers y ejecutar \texttt{git pull}; se carga automáticamente en la siguiente sesión nueva.
    \item La versión actual admite \textbf{descubrimiento autónomo}: al escribir \texttt{/skills} se listan todos los Skills disponibles.
\end{itemize}

\subsection{Resumen de la sección}

Los Skills, mediante Prompt Injection, exponen su directorio en el tercer mensaje User de Codex; el Agent empareja dinámicamente según el query del usuario y usa el comando \texttt{cat} para cargar el \texttt{skill.md} correspondiente, inyectando así el proceso de ingeniería de software en el contexto del modelo.
\section{Skill maestra: Using Superpowers}

\subsection{Flujo maestro}

\texttt{Using Superpowers} es la primera Skill que se activa («use when starting any conversation»). Define la cadena de desarrollo completa:

\begin{enumerate}
    \item \textbf{Brainstorming}(lluvia de ideas): varias rondas de conversación con el usuario para confirmar los límites de la modificación.
    \item \textbf{Writing Plans}(redacción del plan): convierte el diseño en pasos ejecutables.
    \item \textbf{Development}(desarrollo): se escribe primero la prueba y después la implementación, con TDD.
    \item \textbf{Verification}(verificación): antes de dar por terminado, se ejecutan de verdad los comandos para comprobar el resultado.
\end{enumerate}

\begin{importantbox}{Instrucción de activación obligatoria en Skill.md}
El \texttt{skill.md} de Superpowers contiene una instrucción obligatoria envuelta en etiquetas XML: \textit{«sumamente importante: si crees que hay 1\% de probabilidad de poder usar alguna Skill, actívala sin falta».}
\end{importantbox}

\subsection{Definición formal del flujo}

Dentro de la Skill se usa \textbf{el lenguaje Dot}(Graphviz)para definir el diagrama de flujo: describir el SOP en un lenguaje de código formal facilita que el modelo de lenguaje lo asimile. El flujo además exige que el Agente \textbf{declare}(declare)qué Skill está usando y con qué propósito, y define un conjunto de \textbf{Checklist} para verificar si se cumplieron los requisitos.

\subsection{Resumen de la sección}

Using Superpowers es el «punto de entrada al sistema operativo» del Agente: restringe su comportamiento mediante un grafo dirigido definido en el lenguaje Dot, y garantiza, mediante el mecanismo de declaración y el Checklist, que cada paso quede documentado.

\section{Brainstorming(tormenta de ideas)}

\subsection{Objetivos y principios}

Brainstorming resuelve el \textbf{What \& Why}: «qué vamos a construir y por qué lo construimos así».

\begin{itemize}
    \item Se hace una sola pregunta de confirmación a la vez.
    \item Se entiende la intención real del usuario (Purpose), sus restricciones y los criterios de éxito.
    \item Primero se explora el contexto y se proponen \textbf{2--3 opciones}, analizando el Trade-off de cada una.
    \item Una vez que el usuario confirma, se forma y se guarda el \textbf{Design Doc} (\texttt{docs/plans/<date>-<topic>-design.md}).
\end{itemize}

\begin{warningbox}{Prohibiciones en la etapa de Brainstorming}
En la etapa de Brainstorming, \textbf{está prohibido invocar cualquier herramienta de implementación y prohibido escribir código}. El producto de esta etapa es puramente conceptual y de documentos de decisión, centrado en la arquitectura general, el flujo de datos, la interacción entre componentes y el manejo de excepciones.
\end{warningbox}

\subsection{Pasos del proceso}

\begin{enumerate}
    \item Explorar el contexto del proyecto.
    \item Plantear al usuario preguntas de confirmación para delimitar el problema con claridad.
    \item Proponer 2--3 soluciones y analizar el Trade-off.
    \item Presentar el Design Section y esperar la aprobación del usuario.
    \item Una vez aprobado por el usuario, invocar automáticamente el Skill \texttt{Writing Plans}.
\end{enumerate}

\subsection{Resumen de la sección}

El Brainstorming es la etapa obligatoria de esclarecimiento de requisitos; produce el Design Doc, que sienta la base arquitectónica para el plan de implementación posterior. Su disciplina central es «solo diseñar, no programar».


\section{Writing Plans (planes de implementación)}

\subsection{Objetivo y principios}

Writing Plans resuelve el \textbf{How}, es decir, «cómo implementarlo». Convierte el diseño macroscópico en pasos de grano fino \textbf{ejecutables sin contexto}.

\begin{importantbox}{Tres principios fundamentales}
\begin{enumerate}
    \item \textbf{DRY} (Don't Repeat Yourself): evitar escribir código y lógica repetidos.
    \item \textbf{YAGNI} (You Aren't Gonna Need It): no construir funcionalidad que «tal vez sirva después».
    \item \textbf{TDD}: escribir primero las pruebas y después la implementación, con confirmaciones frecuentes e iteraciones en pasos pequeños.
\end{enumerate}
\end{importantbox}

\subsection{Entregable: Plan.md}

El documento del plan (\texttt{Plan.md}) contiene detalles extremadamente específicos:
\begin{itemize}
    \item Qué archivo crear, con su ruta exacta.
    \item Qué líneas modificar.
    \item Qué prueba fallida escribir.
    \item Qué comando exacto ejecutar para verificar.
\end{itemize}

\begin{knowledgebox}{Eliminar la dependencia del contexto}
El estándar para escribir el plan es: suponer que la persona ingeniera (u otro Subagent) que finalmente escribirá el código no sabe nada del proyecto y puede completar el desarrollo basándose únicamente en este plan. Por eso se puede abrir otra sesión y dejar que Codex, Claude Code o una persona ingeniera lo ejecute directamente.
\end{knowledgebox}

\subsection{Opciones tras terminar el plan}

Al terminar de escribir el plan no se escribe código de inmediato, sino que se le pregunta al usuario:
\begin{itemize}
    \item Ejecutarlo de forma incremental en la sesión actual (Sequential Development).
    \item Abrir una sesión nueva para ejecutarlo por lotes (Parallel Agents).
\end{itemize}

\subsection{Resumen de la sección}

Writing Plans «traduce» el diseño arquitectónico en una lista de pasos microscópicos que cualquier persona o Agent puede ejecutar; es el puente clave entre «pensar» y «hacer».

\section{TDD: rojo-verde-refactorización}

\subsection{Disciplina de tres etapas}

\begin{importantbox}{Rojo-verde-refactorización}
\begin{enumerate}
    \item \textbf{Rojo (Red)}: primero se escribe una prueba que \textbf{fallará}; como la funcionalidad todavía no está implementada, la prueba necesariamente marca error.
    \item \textbf{Verde (Green)}: se escribe la \textbf{cantidad mínima} de código para que la prueba pase.
    \item \textbf{Refactorización (Refactor)}: se optimiza la estructura del código manteniéndolo en verde.
\end{enumerate}
\end{importantbox}

\subsection{¿Por qué es indispensable "ver con los propios ojos que la prueba falla"?}

\begin{warningbox}{Una prueba escrita después no demuestra nada}
\begin{itemize}
    \item Después de escribir la prueba hay que ejecutarla y ver el \textbf{error esperado}.
    \item Si la prueba pasa de inmediato, significa que lo que se está probando está mal, o que se está probando una conducta que ya existía.
    \item Escribir primero el código y agregar la prueba después casi siempre solo prueba "la conducta actual del código", no "la conducta que el código debería tener"; de ahí viene la \textbf{deuda técnica}.
\end{itemize}
\end{warningbox}

\subsection{La verificación necesita evidencia}

Antes de declarar que algo está terminado o que una prueba pasó, hay que \textbf{ejecutar de verdad el comando o la herramienta} y revisar la salida, usando abundantes aserciones (assert) con una expectativa clara sobre el resultado.

\subsection{Resumen de la sección}

La disciplina central de TDD es "primero rojo, luego verde, después refactorización": obliga al Agent a usar en cada paso resultados reales de pruebas como evidencia, y elimina el "fingir que algo está terminado".

\section{Modos de ejecución}

\subsection{Sequential Development(desarrollo progresivo)}

En la sesión actual se ejecuta paso a paso según el plan: al terminar cada task, se realiza implementación$\rightarrow$verificación$\rightarrow$Code Review, y después se pasa a la siguiente. Es adecuado para tareas complejas y de alto riesgo que requieren supervisión humana en cada paso.

\subsection{Parallel Agents(Subagent concurrentes)}

Cuando las tareas del plan son independientes entre sí, se despacha un nuevo Subagent para cada tarea y se ejecutan de forma concurrente. Al terminar, se realiza una doble revisión automática (especificación + calidad).

\subsection{Aislamiento con Worktree}

Se usa Git Worktree para aislar la rama actual, y en la ejecución de un solo plan se enfatiza el orden «implementación$\rightarrow$especificación$\rightarrow$Review de la calidad del código».

\subsection{Resumen de la sección}

Superpowers ofrece múltiples modos de ejecución, desde el progresivo de un solo hilo hasta la concurrencia de múltiples Agent, adaptándose con flexibilidad a escenarios de desarrollo de distinta complejidad.

\section{Systematic Debugging(depuración sistemática)}

\subsection{Proceso de cuatro fases}

Se activa al encontrar un error o una prueba fallida, \textbf{está prohibido adivinar a ciegas o corregir por ensayo y error}:

\begin{enumerate}
    \item \textbf{Investigación de la causa raíz}(Root Cause Analysis): encontrar la causa real.
    \item \textbf{Análisis de patrones}(Pattern Analysis): identificar patrones de problemas que se repiten.
    \item \textbf{Plantear una hipótesis}(Hypothesis): proponer una corrección basada en evidencia.
    \item \textbf{Implementación mínima y verificación}(Minimal Fix \& Verify).
\end{enumerate}

\begin{warningbox}{Si tres ciclos no lo resuelven, hay un problema de arquitectura}
Si el ciclo de cuatro fases anterior se repite tres veces sin resolver el problema, esto indica que probablemente existe un problema a nivel de arquitectura y es necesario volver al Brainstorming para reconsiderar el diseño.
\end{warningbox}

\subsection{Skill de verificación antes de dar por terminado}

A la IA le gusta «fingir que terminó» o mostrar exceso de confianza. Este Skill exige: antes de anunciar que una tarea está completa, que las pruebas pasaron o que está lista para enviarse, es obligatorio \textbf{ejecutar realmente el comando de verificación} y declarar la finalización con base en la salida real.

\subsection{Resumen de la sección}

La depuración sistemática sustituye la corrección «a capricho» por un ciclo cerrado de cuatro fases; si tres intentos no bastan, el problema se eleva a un problema de arquitectura, para evitar hundirse cada vez más en una dirección equivocada.
\section{Sistema de Code Review}

\subsection{Momento y principios del Review}

\begin{importantbox}{{Review Early, Review Often}}
\begin{itemize}
    \item No se revisa solo antes de fusionar con la rama principal, sino que \textbf{después de completar cada subtarea de grano fino} se debe activar la revisión de manera obligatoria.
    \item Evita que un error de diseño «contamine» la siguiente tarea.
    \item Entrada del Code Review: los requisitos/el plan original del usuario + el Git Diff.
\end{itemize}
\end{importantbox}

\subsection{Checklist para solicitar un Code Review}

\begin{enumerate}
    \item \textbf{Calidad del código}: separación de responsabilidades, desacoplamiento, manejo de errores, código duplicado, cobertura de casos límite.
    \item \textbf{Arquitectura}: si los roles de diseño son razonables, escalabilidad, impacto en el rendimiento, riesgos de seguridad.
    \item \textbf{Pruebas}: si se prueba la lógica real (y no solo con Mock), cobertura de casos límite, pruebas de integración.
    \item \textbf{Requisitos}: si se cumplen todos los requisitos del plan, sin proliferación de funcionalidades, y se alcanza un nivel de acabado apto para producción.
\end{enumerate}

\begin{warningbox}{Líneas rojas de conducta}
\begin{itemize}
    \item Está terminantemente prohibido decir «se ve bien» sin haber revisado el código.
    \item Está terminantemente prohibido marcar un problema secundario como un error fatal.
    \item Está terminantemente prohibido ser ambiguo — se debe llegar a una conclusión final.
\end{itemize}
\end{warningbox}

\subsection{Criterios para recibir un Code Review}

La esencia del Code Review es \textbf{una evaluación técnica basada en el código objetivo}, no una actuación social:

\begin{itemize}
    \item \textbf{Buscar el rigor técnico}: las decisiones se basan en hechos, resultados de pruebas y el estado actual de la arquitectura del sistema, no en la intuición o la autoridad.
    \item \textbf{Rechazar el acuerdo performativo}: no aceptar a ciegas las sugerencias del Reviewer sin haber verificado antes su viabilidad técnica.
    \item \textbf{Principio YAGNI}: si el Reviewer sugiere añadir una funcionalidad que el sistema actual no invoca en absoluto, debe rechazarse.
    \item \textbf{Verificar antes de actuar}: al recibir la retroalimentación, el primer paso es verificarla en la base de código, no modificar de inmediato.
    \item \textbf{Objeción técnica bien fundamentada} (Push Back): cuando se detecta que la sugerencia del Reviewer es errónea, se refuta con razonamiento técnico — la objeción no busca discutir, sino proteger la base de código.
\end{itemize}

\begin{knowledgebox}{Eliminar la ambigüedad}
Si el Reviewer da 6 observaciones y 2 de ellas no quedan claras, \textbf{está terminantemente prohibido hacer primero los 4 puntos que sí se entienden y dejar los 2 que no se entienden para el final}. Hay que detenerse de inmediato y preguntar primero todos los detalles con claridad, porque los módulos de software están altamente acoplados entre sí, y una comprensión parcial llevaría a una dirección de implementación equivocada, con un costo de retrabajo enorme.
\end{knowledgebox}

\begin{importantbox}{Los hechos hablan más que las palabras}
Cuando la observación del Reviewer es correcta, está terminantemente prohibido usar frases vacías de agradecimiento o elogio. Lo correcto es \textbf{corregirlo directamente y exponer los hechos de forma breve}. En una cultura de ingeniería eficiente, el código mismo es la mejor respuesta.
\end{importantbox}

\subsection{SOLID y la separación de responsabilidades}

Patrones de diseño arquitectónico que intervienen en el Code Review:

\begin{enumerate}
    \item \textbf{Principio de responsabilidad única} (SRP): evita que una clase o un archivo asuma todas las tareas; se debe desacoplar.
    \item \textbf{Principio de abierto/cerrado} (OCP): abierto a la extensión, cerrado a la modificación.
    \item \textbf{Principio de sustitución de Liskov} (LSP): una subclase puede sustituir a su clase padre.
    \item \textbf{Principio de segregación de interfaces} (ISP): las interfaces deben ser de grano fino.
    \item \textbf{Principio de inversión de dependencias} (DIP): depender de abstracciones y no de implementaciones concretas.
\end{enumerate}

La \textbf{separación de responsabilidades} es una filosofía arquitectónica de alto nivel: divide el software en módulos independientes con responsabilidades distintas, de modo que cada módulo resuelve solo un aspecto específico.

\subsection{Resumen de la sección}

El sistema de Code Review comprende dos Skills, «solicitar la revisión» y «recibir la revisión»; mediante el Checklist y las líneas rojas de conducta se garantiza el rigor técnico de la revisión, mientras que los principios SOLID y la separación de responsabilidades guían la evaluación de la arquitectura.

\section{Reutilización de código y DRY}

En escenarios de Vibe Coding, el modelo está limitado por la ventana de contexto, por lo que tiende a modificar el código habiendo visto solo una parte de él, lo que provoca que las funciones de utilidad (Utils) acumulen una gran cantidad de código redundante y duplicado.

Superpowers insiste una y otra vez, tanto en \texttt{Writing Plans} como en la revisión de código:

\begin{itemize}
    \item \textbf{DRY} (Don't Repeat Yourself): eliminar el código duplicado.
    \item \textbf{YAGNI} (You Aren't Gonna Need It): eliminar de manera decidida el diseño excesivo y la abstracción excesiva.
\end{itemize}

Además, Superpowers ofrece el Skill \texttt{Writing Skills}, que insiste en depurar el proceso de resolución de problemas hasta convertirlo en un \textbf{patrón reutilizable}, en lugar de un script de un solo uso.

\subsection{Resumen de la sección}

DRY y YAGNI son los dos principios que atraviesan todo Superpowers: el primero evita la redundancia y el segundo evita la sobreingeniería; juntos constituyen el equilibrio dinámico de la calidad del código.

\section{Cómo escribir buenos Skills}

El expositor resumió los elementos para escribir Skills de alta calidad:

\begin{enumerate}
    \item \textbf{Suficientemente abstractos}: los Skills deben ser generales (Common), estar basados en principios (Principled) y ser reutilizables.
    \item \textbf{Descripción formalizada}: usar un lenguaje de descripción muy Formal y lógico (como diagramas de flujo en lenguaje Dot).
    \item \textbf{Condensación de la experiencia experta}: deben ser la síntesis de la práctica de desarrolladores de software con experiencia.
    \item \textbf{Capacidad de Writing}: se requiere una buena capacidad de escritura técnica, aunque se puede recurrir a un modelo de lenguaje para pulirla.
\end{enumerate}

\subsection{Resumen de la sección}

Un buen Skill es la cristalización de «experiencia experta + expresión formalizada»: requiere tanto una comprensión profunda de la ingeniería como una descripción estructurada que el modelo pueda ejecutar con facilidad.


\section{Síntesis y ampliación}

Esta sesión tomó como eje el proyecto Superpowers para explicar de manera sistemática cómo «conducir» (Harness) al AI Coding Agent en la era del Vibe Coding, de modo que escriba código «correcto, bueno y preciso». Repaso de los puntos clave:

\begin{enumerate}
    \item \textbf{Mecanismo de carga de Skills}: mediante Prompt Injection se expone el directorio de Skills en el mensaje de sistema de Codex, y el Agent los relaciona y carga de manera dinámica.
    \item \textbf{Flujo de control general}: Using Superpowers $\rightarrow$ Brainstorming $\rightarrow$ Writing Plans $\rightarrow$ TDD Development $\rightarrow$ Verification.
    \item \textbf{Disciplina de TDD}: las tres etapas de rojo-verde-refactorización, donde «ver con los propios ojos que la prueba falla» es el primer paso que no puede omitirse.
    \item \textbf{Depuración sistemática}: un ciclo cerrado de cuatro etapas; si tres intentos no dan resultado, el problema se escala a un problema de arquitectura.
    \item \textbf{Sistema bidireccional de Code Review}: solicitar la revisión (impulsada por un Checklist) + recibir la revisión (rigor técnico + Push Back).
    \item \textbf{DRY \& YAGNI}: los dos principios que atraviesan todo el proceso.
\end{enumerate}

El significado de Superpowers no es solo un conjunto de Prompts, sino «codificar de manera rígida» un conjunto verificado de las mejores prácticas de ingeniería de software como el sistema operativo normativo de la era de la IA. Representa la \textbf{Harness Engineering}: mediante el Skill como esa «rienda», conduce a un Coding Agent cada vez más capaz y encuentra el equilibrio óptimo entre eficiencia y calidad.

\subsection{Lecturas adicionales}

\begin{itemize}
    \item Repositorio de GitHub de Superpowers
    \item Kent Beck, \textit{Test-Driven Development: By Example}
    \item Robert C. Martin, \textit{Clean Code}
    \item Martin Fowler, \textit{Refactoring: Improving the Design of Existing Code}
\end{itemize}

%% --- Fin del contenido --- %%

\end{document}
